\documentclass[11pt]{article}
\usepackage[margin=0.88in]{geometry}
\usepackage[T1]{fontenc}
\usepackage[utf8]{inputenc}
\usepackage{amsmath,amssymb,amsthm}
\usepackage{graphicx}
\usepackage{xcolor}
\usepackage{microtype}
\usepackage[colorlinks=true,urlcolor=blue!55!black,citecolor=blue!55!black]{hyperref}

\newtheorem{theorem}{Theorem}
\newtheorem{lemma}[theorem]{Lemma}
\newcommand{\D}{\mathbb D}
\newcommand{\C}{\mathbb C}
\newcommand{\T}{\mathbb T}
\newcommand{\Acal}{\mathcal A}
\newcommand{\Log}{\operatorname{Log}}

\title{A logarithmic strip tail destroys every finite-order\\Laurent decomposition at an internal tangency}
\author{Candidate full negative answer to Question 34 of arXiv:1605.08289}
\date{11 August 2026}

\begin{document}
\maketitle

\begin{center}
\fbox{\parbox{0.92\textwidth}{\textbf{Status: candidate counterexample, likely valid; full for every finite $p$ in the source's parameter set.}
For each integer $p\geq0$ and every round disc internally tangent to the unit
disc, an explicit $F_p\in A^p(\Omega)$ is constructed which is not in
$A^p(\D)+A^p_0(\widehat\C\setminus\overline{D'})$.  The construction turns
the cusp into a vertical strip and leaves a one-sided logarithmic tail in the
$p$-th derivative.}}
\end{center}

\section*{Source question}

Let $D'$ be a disc contained in $\D$ whose boundary is tangent to $\T$ at
$1$, and put $\Omega=\D\setminus\overline{D'}$.  The source proves the
Laurent decomposition for $p=\infty$ and then asks:
\begin{quote}
For $p<+\infty$, is there an $F\in A^p(\Omega)$ that cannot be written as
the sum of a function in $A^p(\D)$ and another in
$A^p_0(\widehat\C\setminus\overline{D'})$?
\end{quote}
This is Question~34 on PDF page 13 of Georgakopoulos's paper
\cite[Question~34]{source}.  The page containing the question and its Cauchy
transform criterion is reproduced below.

\begin{center}
\includegraphics[width=0.91\textwidth]{figures/open_problem_crop.png}
\end{center}

Here $p\in\{0,1,2,\ldots\}$, as in the source's definition of $A^p$.

\section*{Proof intuition}

Inversion at the tangency point, $w=(1-z)^{-1}$, sends the crescent $\Omega$
to a vertical strip.  Its two infinite ends both represent the same boundary
point $z=1$.  An analytic cutoff on the strip can tend to $1$ at one end and
to $0$ exponentially at the other, while the resulting function still has a
single limit at the tangency point.

We multiply that cutoff by $w^{-p}/\Log(2-iw)$.  The factor $w^{-p}$ absorbs
the $p$ differentiations in the $z$ variable; after $p$ derivatives the
positive end still has size $1/\log|w|$, whereas the negative end is
exponentially small.  On the outer circle this becomes a one-sided boundary
tail $1/\log(1/t)$.  Its Cauchy projection grows like
$\log\log(1/(1-\rho))$ as $\rho\uparrow1$.  If an $A^p$ Laurent decomposition
existed, differentiating it $p$ times would make this projection an element
of $A(\D)$, which must have a finite limit at $1$.

\section*{The counterexample}

Write
\[
 D'=D'_r:=\{z:|z-(1-r)|<r\},\qquad 0<r<1,
\]
which is the general disc internally tangent to $\T$ at $1$.  Set
\[
 a=\frac12,\qquad b=\frac1{2r},\qquad L=b-a,
 \qquad W(z)=\frac1{1-z}.
\]
For $a<\Re w<b$, define
\begin{align*}
 E(w)&=\exp\!\left(\frac{\pi i(w-a)}{L}\right),
 &\chi(w)&=\frac1{1-iE(w)},\\
 \Lambda(w)&=\Log(2-iw),
 &\Phi_p(w)&=\frac{w^{-p}\chi(w)}{\Lambda(w)},
\end{align*}
where $\Log$ is the principal logarithm.  Our function is
\begin{equation}\label{eq:Fp}
 F_p(z)=\Phi_p(W(z)),\qquad z\in\Omega.
\end{equation}

\begin{theorem}\label{thm:main}
For every integer $p\geq0$, the function $F_p$ in \eqref{eq:Fp} belongs to
$A^p(\Omega)$ but
\[
 F_p\notin A^p(\D)+A^p_0(\widehat\C\setminus\overline{D'_r}).
\]
Consequently Question~34 has an affirmative answer for every finite $p$.
\end{theorem}

\section*{Proof}

\subsection*{1. The strip and its two ends}

For $w=W(z)$ one has
\[
 |z|<1\iff \Re w>\frac12,
 \qquad
 |z-(1-r)|>r\iff \Re w<\frac1{2r}.
\]
Thus $W$ maps $\Omega$ conformally onto the strip
$S=\{a<\Re w<b\}$.  The two limits $\Im w\to+\infty$ and
$\Im w\to-\infty$ both correspond to $z\to1$ through the two arms of the
cusp.

The map $E$ sends $S$ to the upper half-plane.  Uniformly for
$a\leq\Re w\leq b$,
\begin{equation}\label{eq:chiends}
 \chi(w)=1+O(e^{-\pi\Im w/L})\quad(\Im w\to+\infty),
 \qquad
 \chi(w)=O(e^{-\pi|\Im w|/L})\quad(\Im w\to-\infty),
\end{equation}
and the same estimates, up to polynomial factors, hold after any fixed
number of $w$-derivatives.  Indeed, the only zero of $1-i\zeta$ is
$\zeta=-i$, outside the closed upper half-plane.

Moreover $2-iw$ lies in the open lower half-plane throughout the closed
strip, because $\Im(2-iw)=-\Re w<0$.  Hence $\Lambda$ is holomorphic on a
neighborhood of every finite point of the closed strip and never vanishes
there.  At either end,
\begin{equation}\label{eq:Lambda}
 |\Lambda(w)|\sim\log|w|.
\end{equation}

\subsection*{2. Finite-order regularity at the cusp}

Since $dW/dz=W^2$, differentiation in $z$ becomes
\[
 \mathcal T:=w^2\frac{d}{dw}.
\]
Repeated product differentiation using \eqref{eq:chiends}--\eqref{eq:Lambda}
gives, for $0\leq k\leq p$ and $\Im w\to+\infty$,
\begin{equation}\label{eq:Tasymp}
 \mathcal T^k\Phi_p(w)
 =(-1)^k\frac{p!}{(p-k)!}\frac{w^{k-p}}{\Lambda(w)}
 +O\!\left(\frac{|w|^{k-p}}{|\Lambda(w)|^2}\right)+o(1),
\end{equation}
uniformly on the closed strip.  Terms in which a derivative hits $\chi$
are exponentially small; a derivative hitting $\Lambda^{-1}$ gains one
power of $w$ but also one additional inverse power of $\Lambda$.  This also
proves \eqref{eq:Tasymp} by induction.  At the negative end,
\begin{equation}\label{eq:Tminus}
 \mathcal T^k\Phi_p(w)=O(e^{-c|\Im w|}|w|^{N_k})=o(1)
\end{equation}
for suitable $c>0$ and $N_k$.

Equations \eqref{eq:Tasymp} and \eqref{eq:Tminus} show that
$\mathcal T^k\Phi_p(w)\to0$ at both ends for every $k\leq p$ (also when
$k=p$, because $|\Lambda(w)|\to\infty$).  Away from the cusp all functions
are holomorphic on a neighborhood of the two boundary arcs.  Therefore each
$F_p^{(k)}$, $0\leq k\leq p$, extends continuously to $\overline\Omega$ by
the value $0$ at $1$.  This proves
\begin{equation}\label{eq:Ap}
 F_p\in A^p(\Omega).
\end{equation}

For later use, put $Q_p=F_p^{(p)}$ and $A_p=(-1)^p p!$.  The case $k=p$ of
\eqref{eq:Tasymp} says
\begin{equation}\label{eq:Qstrip}
 Q_p(z)=\frac{A_p}{\Lambda(W(z))}
 +O\!\left(\frac1{|\Lambda(W(z))|^2}\right)
 \quad(\Im W(z)\to+\infty),
\end{equation}
while $Q_p$ is exponentially small, up to a polynomial factor, at the other
end.

\subsection*{3. The outer Cauchy projection diverges}

On the outer circle, with $z=e^{it}$,
\begin{equation}\label{eq:wt}
 W(e^{it})=\frac12+\frac{i}{2}\cot\frac t2.
\end{equation}
Thus \eqref{eq:Qstrip} yields, as $t\downarrow0$,
\begin{equation}\label{eq:Qplus}
 \frac{Q_p(e^{it})}{A_p}
 =\frac1{\log(1/t)}+O\!\left(\frac1{\log^2(1/t)}\right),
\end{equation}
whereas for $t\uparrow0$ from below,
\begin{equation}\label{eq:Qminus}
 Q_p(e^{it})=O\!\left(|t|^{-N}e^{-c/|t|}\right)
\end{equation}
for some $c,N>0$.

For $0<\rho<1$, let
\[
 \mathcal C_p(\rho)=\frac1{2\pi i}\int_{\T}
 \frac{Q_p(\zeta)}{\zeta-\rho}\,d\zeta
 =\frac1{2\pi}\int_{-\pi}^{\pi}Q_p(e^{it})
 \frac{e^{it}}{e^{it}-\rho}\,dt.
\]
The exact imaginary part of the kernel is
\begin{equation}\label{eq:kernel}
 \Im\frac{e^{it}}{e^{it}-\rho}
 =-\frac{\rho\sin t}{(1-\rho)^2+2\rho(1-\cos t)}.
\end{equation}
Put $\delta=1-\rho$.  On $4\delta\leq t\leq\varepsilon$, the right side of
\eqref{eq:kernel} is $-1/t$ up to an error whose integral against
$1/\log(1/t)$ remains bounded.  The real part of the kernel, multiplied by
the imaginary error in \eqref{eq:Qplus}, also has uniformly bounded integral.
The interval $0<t<4\delta$ contributes $O(1/\log(1/\delta))$; the negative
arc is uniformly bounded by \eqref{eq:Qminus}; and arcs away from $1$ are
uniformly bounded.  Consequently
\begin{align}
 \Im\frac{\mathcal C_p(\rho)}{A_p}
 &=-\frac1{2\pi}\int_{4\delta}^{\varepsilon}
   \frac{dt}{t\log(1/t)}+O(1)\notag\\
 &=-\frac1{2\pi}\log\log\frac1\delta+O(1)
 \longrightarrow-\infty.\label{eq:loglog}
\end{align}
In particular, the outer Cauchy projection of $Q_p$ has no finite limit at
$1$ from inside $\D$.

\subsection*{4. Contradiction to a Laurent decomposition}

Suppose that $F_p=g+h$ with
\[
 g\in A^p(\D),\qquad
 h\in A^p_0(\widehat\C\setminus\overline{D'_r}).
\]
Differentiating $p$ times gives $Q_p=g^{(p)}+h^{(p)}$, with
$g^{(p)}\in A(\D)$.  Since $h^{(p)}$ is holomorphic on the exterior of the
unit circle and vanishes at infinity, its Cauchy integral over $\T$ is zero
inside $\D$.  Cauchy's formula therefore gives
\[
 \mathcal C_p(z)=g^{(p)}(z),\qquad z\in\D.
\]
The right side extends continuously to $z=1$, contradicting
\eqref{eq:loglog}.  This proves Theorem~\ref{thm:main}. \qed

\section*{Clarification of the source's final sentence}

The source later asks, ``in particular,'' whether some $F\in A(\D)$ can have
a divergent inner limit of its outer-circle Cauchy transform.  Literally this
subquestion has the immediate negative answer
\[
 C_{\T}(F)(z)=F(z)\longrightarrow F(1),\qquad z\to1, z\in\D,
\]
by Cauchy's formula and the definition of the disc algebra.  The substantive
Question~34 concerns $F\in A^p(\Omega)$; Theorem~\ref{thm:main} answers that
question fully.  For $p>0$ the proof uses the Cauchy projection of the
$p$-th derivative, rather than claiming that the projection of $F_p$ itself
diverges.

\section*{Verification and novelty bounds}

The strip identities, the two endpoint asymptotics, and the Cauchy-kernel
estimate are symbolic.  The included script
\texttt{code/check\_asymptotics.py} independently differentiates the formula
for $p=0,1,2,3$ and checks the one-sided boundary ratio and the predicted
log--log growth numerically; it is a sanity check, not part of the proof.

A bounded novelty search on 11 August 2026 covered the run's registry,
solution, attempt, and proof-gap indexes; the exact arXiv id and title;
``Question 34'' with Laurent/$A^p$ keywords; the author name; and
``internally tangent circles'' with the decomposition and Cauchy-transform
keywords.  It found the source paper and later work on $A^p$ boundary traces
for Jordan domains, but no later answer to this tangent-disc question or this
strip/logarithmic construction.  Novelty confidence is therefore moderate,
not definitive.

Human review should focus on the induction in \eqref{eq:Tasymp}, uniformity at
both strip ends, and the bounded-error reduction from \eqref{eq:kernel} to
\eqref{eq:loglog}.  No numerical statement is used in the proof.

\enlargethispage{4\baselineskip}
\begin{thebibliography}{1}
\bibitem{source}
N.~Georgakopoulos,
\emph{Extensions of the Laurent Decomposition and the spaces $A^p(\Omega)$},
arXiv:1605.08289 (2016).
\end{thebibliography}

\end{document}
